\begin{studiumworked}{M12. Superficie y garganta variables: modelo normalizado · Intermedio}
\textbf{Enunciado.} En un modelo académico de regresión se conserva densidad, temperatura y $c^*$. La ley es $r_b/r_{ref}=(p_c/p_{ref})^n$, con $n=0.344$. Entre dos estados el área activa se multiplica por $1.24$ y la garganta por 1.02. Calcule razones de presión, regresión y caudal. Determine la razón de empuje si $C_F$ se mantiene.

\textbf{Datos.} Exponente local: $0.344$; Razón de área activa: $1.24$; Razón de garganta: $1.02$.

\textbf{Incógnitas.} Cocientes de presión, regresión, masa y empuje.

\textbf{Hipótesis.} Cuasiestacionario; descarga ahogada. Propiedades y CF fijos; ejemplo normalizado sin asignación a propulsante real.

\textbf{Modelo físico.} Equilibrio de generación y descarga.
\[\rho_pA_br_{ref}(p_c/p_{ref})^n=p_cA_t/c^*\]
\[P^{1-n}=\frac{A_{b,2}/A_{b,1}}{A_{t,2}/A_{t,1}}\]
\[P=(\alpha/\tau)^{1/(1-n)},\quad r_{b,2}/r_{b,1}=P^n\]
\[\dot m_2/\dot m_1=P\tau,\quad F_2/F_1=P\tau\ (C_F\ fijo)\]
\textbf{Desarrollo.}
\begin{enumerate}
\item Dividir los equilibrios elimina los parámetros dimensionales comunes. Queda $P^{1-0.344}=1.24/1.02$. Se despeja elevando a $1/(1-n)$, porque $n<1$.
\item Sustituir da $P=1.346794$. La regresión aumenta por un factor $P^n=1.107846$. El caudal generado cambia por $\alpha P^n=1.37373$.
\item La descarga produce independientemente $P\tau=1.37373$, igual al caudal generado. Con $C_F$ fijo, $F=C_Fp_cA_t$ da el mismo factor de empuje; usar sólo $P$ omitiría el cambio de garganta.
\end{enumerate}
\textbf{Resultados.} $p_2/p_1=1.346794$; $r_{b,2}/r_{b,1}=1.107846$; $\dot m_2/\dot m_1=F_2/F_1=1.37373$.

\textbf{Interpretación.} El aumento de garganta compensa parte del efecto de área activa. La hipótesis de CF fijo es adicional y podría fallar si cambia la expansión o presión ambiente.

\textbf{Comprobación.} El cierre $\alpha P^n=P\tau$ se verifica por ambas rutas. Todos los cocientes son adimensionales. Para $\alpha=\tau=1$ se recupera equilibrio inicial.

\textbf{Límites.} No se incluye almacenamiento, erosión predictiva, cinética de superficie ni inestabilidad acústica. Un exponente local no es una garantía de estabilidad real.
\end{studiumworked}
